\chapter{Exterior algebra, tableaux, and Schur constructions}
\label{chap:exterior-schur}

\section{The algebra needed to describe a space of complexes}

A matrix equation such as \(BA=0\) has several simultaneous meanings.
It states that two linear maps compose to zero; it gives polynomial
equations in their entries; and it makes every exterior power of their
composition vanish.  The first interpretation controls homology.  The
second defines an algebraic scheme.  The third produces the determinant
identities used to describe its coordinate ring.  Standard monomial
theory makes the third interpretation effective: it chooses products of
minors that form a basis, and supplies relations that express all other
products in that basis.

This requires more than remembering that determinants are alternating.
One must keep track of integral coefficients, the direction of each
linear map, the dual action on coordinate functions, and the distinction
between a basis of a module and multiplication in an algebra.  These
chapters develop that bookkeeping explicitly.  The elementary arguments
are proved here.  The general standard basis and normality theorems are
identified as external results, with their hypotheses and their precise
role stated.

Throughout this chapter \(R\) is a commutative ring with identity and
\(V\) is a finite free \(R\)-module.  Statements involving irreducible
representations will instead specify an algebraically closed field of
characteristic zero.  Maintaining this distinction matters for orders:
an argument using division by a hook length cannot silently be
transported to an order over a residue field whose characteristic divides
that hook length.

\section{Exterior powers and determinant coordinates}

\begin{definition}
The exterior algebra \(\bigwedge V\) is the quotient of the tensor algebra
of \(V\) by the two-sided ideal generated by \(v\otimes v\), for all
\(v\in V\).  Its degree \(p\) part is \(\bigwedge^p V\).
We use \(\bigwedge^0V=R\).  For an ordered basis
\(e_1,\ldots,e_n\), write
\[
 e_I=e_{i_1}\wedge\cdots\wedge e_{i_p},
 \qquad I=\{i_1<\cdots<i_p\}\subseteq\{1,\ldots,n\}.
\]
\end{definition}

The relation \(v\wedge v=0\) is part of the definition.  In characteristic
two it does not follow merely from the formula
\(v\wedge w=-w\wedge v\).  Expanding \((v+w)\wedge(v+w)=0\) does,
however, prove the latter formula over every \(R\).

\begin{proposition}[The exterior basis]
The elements \(e_I\), with \(|I|=p\), form an \(R\)-basis of
\(\bigwedge^pV\).  In particular its rank is \(\binom np\), interpreted
as zero when \(p>n\).
\end{proposition}
\begin{proof}
Multilinearity expands any pure wedge in the basis vectors.  A term
with a repeated index vanishes; adjacent interchanges put the remaining
indices in increasing order, so the displayed elements span.  For every
ordered \(p\)-element set \(J\), define an alternating \(p\)-linear form
on \(V^p\) by taking the determinant of the coordinates in the rows
indexed by \(J\).  The universal property of the exterior quotient
gives a functional \(\ell_J:\bigwedge^pV\to R\).
Its value on \(e_I\) is \(1\) for \(I=J\), and zero otherwise.
Applying these functionals to a proposed linear relation proves
independence, without dividing by any integer.
\end{proof}

A linear map \(T:V\to W\) therefore gives a linear map
\(\bigwedge^pT\) by applying \(T\) to each factor.  If \(T\) has
matrix \(M\), then the coefficient of \(f_I\) in
\((\bigwedge^pT)e_J\) is the minor \(\det M_{I,J}\).
This statement is a direct expansion of the determinant, so it works
over every commutative coefficient ring.

\begin{proposition}[Cauchy--Binet with specified indices]
Let \(A:R^a\to R^b\) and \(B:R^b\to R^c\).  If \(I\subseteq[c]\)
and \(J\subseteq[a]\) have size \(p\), then
\[
 \det(BA)_{I,J}
   =\sum_{\substack{K\subseteq[b]\\ |K|=p}}
      \det B_{I,K}\det A_{K,J}.
\]
\end{proposition}
\begin{proof}
Functoriality on pure wedges gives
\(\bigwedge^p(BA)=(\bigwedge^pB)(\bigwedge^pA)\).
Taking its \(I,J\) matrix entry in the exterior bases gives the formula.
This also proves it when \(p>b\), since the sum is empty and the
intermediate exterior module is zero.
\end{proof}

If \(BA=0\) and \(p>0\), these sums vanish.  An individual summand need
not vanish.  Already for a row \(B=(b_1,b_2)\) and column
\(A=(a_1,a_2)^{\mathsf T}\), the equation is
\(b_1a_1+b_2a_2=0\), not the pair of equations \(b_1a_1=b_2a_2=0\).
Confusing a determinant identity with termwise vanishing destroys the
coordinate ring.  The mixed straightening relations below retain the
necessary sums and signs.

\section{Exterior coproducts and shuffle signs}

For disjoint increasing index sets \(I,J\) whose union is
\(\{1,\ldots,p+q\}\), let \(\epsilon(I,J)\) be the sign of the
permutation that first lists \(I\), then \(J\).

\begin{definition}
The exterior coproduct of bidegree \((p,q)\) is
\[
 \Delta_{p,q}(v_1\wedge\cdots\wedge v_{p+q})
  =\sum_{\substack{I\sqcup J=[p+q]\\ |I|=p}}
    \epsilon(I,J)\,v_I\otimes v_J.
\]
Here the elements in each wedge are taken in their original increasing
index order.
\end{definition}

\begin{proposition}[Shuffle identities over a ring]
The map \(\Delta_{p,q}\) is well defined and natural in \(V\).
Its iterates are coassociative in the graded sense.  If
\(m_{p,q}(u\otimes v)=u\wedge v\), then
\[
 m_{p,q}\Delta_{p,q}=\binom{p+q}{p}\,\mathrm{id}.
\]
\end{proposition}
\begin{proof}
The expression is multilinear.  Interchanging two adjacent input
vectors pairs the summands whose subsets exchange those indices, or
changes the sign inside the corresponding exterior factor.  Thus it
changes the overall sign.  If the two inputs agree, the paired
summands cancel and the other terms contain a repeated wedge.
Consequently the expression is alternating and descends to the
exterior quotient.  Applying a linear map to every input commutes
with every summand, proving naturality.

For coassociativity, expand both ways of splitting a wedge into
three factors of sizes \(p,q,r\).  Both expansions are indexed by
ordered partitions of \([p+q+r]\) into three subsets of those sizes.
In each term, the product of the two shuffle signs is the sign of
the permutation that lists the three subsets in order.  The terms
therefore agree.  Finally, each summand of \(m\Delta\) acquires a
second copy of its shuffle sign when its factors are reordered.
It becomes the original wedge, and there are \(\binom{p+q}{p}\)
summands.
\end{proof}

The final identity is an integral identity, not an instruction to
divide by the binomial coefficient.  For example, dividing
\(\Delta_{1,1}\) by two is unavailable over \(\mathbb F_2\).
Integral Schur constructions use shuffle maps themselves rather than
normalized averaging projectors.

The closely related contraction operation will connect this algebra
to complexes.  For \(\ell\in V^*=\Hom_R(V,R)\), put
\[
 \iota_\ell(v_1\wedge\cdots\wedge v_p)
  =\sum_{j=1}^p(-1)^{j-1}\ell(v_j)
       v_1\wedge\cdots\widehat{v_j}\cdots\wedge v_p.
\]

\begin{lemma}[Contraction and a contracting homotopy]
One has \(\iota_\ell^2=0\).  If \(\ell(u)=1\), the operation
\(h(\omega)=u\wedge\omega\) satisfies
\(\iota_\ell h+h\iota_\ell=\mathrm{id}\).  Thus the contraction
complex is exact.
\end{lemma}
\begin{proof}
In \(\iota_\ell^2\), the term deleting positions \(i<j\) appears
twice.  Deleting \(i\) first changes the later position of \(j\)
by one, so the two signs are opposite.  They cancel over \(R\).
Expanding the definition on \(u\wedge\omega\) gives
\(\iota_\ell(u\wedge\omega)=\ell(u)\omega-u\wedge\iota_\ell\omega\).
This is the asserted homotopy identity.  A cycle \(\omega\)
therefore equals \(\iota_\ell h\omega\), so is a boundary.
\end{proof}

This is the finite algebraic mechanism behind several Koszul
exactness arguments.  In the de Rham symbol complex at a nonzero
covector \(\xi\), choose a vector \(u\) with \(\xi(u)=1\) and exchange
the roles of wedge and contraction to obtain the same homotopy.
The construction proves exactness of that symbol sequence.  It
contains no estimate uniform in a degenerating family of covectors
or connections; such an estimate is a separate analytic issue.

\section{Partitions and a characteristic-free Schur presentation}

A partition \(\lambda=(\lambda_1\ge\cdots\ge\lambda_s>0)\) is
drawn with \(\lambda_i\) boxes in row \(i\).  Its transpose
\(\lambda'\) records column heights.  Write these heights as
\(c_1\ge\cdots\ge c_{\lambda_1}>0\), and set
\[
 E_\lambda(V)=
  \bigwedge^{c_1}V\otimes\cdots\otimes\bigwedge^{c_{\lambda_1}}V.
\]
A tensor in the exterior basis gives a filling of the diagram whose
columns are strictly increasing.  A filling is
\emph{semistandard} here when its columns are strictly increasing
and its rows are weakly increasing.  Some sources transpose this
convention and call the resulting tableaux standard; the convention
is fixed here to prevent a hidden transpose in the Schur index.

For adjacent column heights \(p\ge q\) and \(1\le t\le q\), consider
\[
 \begin{split}
 \theta_{p,q,t}:\ 
 \bigwedge^{p+t}V\otimes\bigwedge^{q-t}V
 &\longrightarrow \bigwedge^pV\otimes\bigwedge^qV,\\
 \theta_{p,q,t}
 &= (\mathrm{id}\otimes m_{t,q-t})
       (\Delta_{p,t}\otimes\mathrm{id}).
 \end{split}
\]
Tensor this map with the identity on the other columns.

\begin{definition}
The Schur module \(S_\lambda(V)\), with the row-partition convention
used here, is the quotient of \(E_\lambda(V)\) by the sum of the
images of these adjacent-column shuffle maps.
\end{definition}

The following statement is the general basis theorem used in this
chapter.  The integral shuffle presentation and its standard basis
are proved in Akin--Buchsbaum--Weyman, Section II.2, with transposed
diagram conventions \cite{smtABW1982}.

\begin{theorem}[External general Schur basis theorem]
For finite free \(V\) with ordered basis, the classes of
semistandard tableaux of shape \(\lambda\) form an \(R\)-basis of
\(S_\lambda(V)\).  This construction commutes with extension of
the coefficient ring.  In particular it is zero if the first
column is longer than \(\rank V\).
\end{theorem}

There are two logically distinct parts to this theorem.  Shuffle
relations move a violating entry toward an earlier column and
produce tableaux lower in a finite tableau order, which proves
spanning.  Independence uses a triangular comparison with monomials
in a tensor product of symmetric powers; the comparison matrix
has diagonal entries one.  The latter fact is what makes the
statement valid over \(\mathbb Z\), rather than only over a field
where a nonzero determinant could be inverted.  The complete
general triangularity argument is the cited result; the following
cases give full proofs without invoking it.

\begin{proposition}[Three elementary Schur cases]
For finite free \(V\),
\[
 S_{(d)}(V)=\Sym^d(V),\qquad
 S_{(1^d)}(V)=\bigwedge^dV.
\]
Moreover the sequence
\[
 0\longrightarrow\bigwedge^3V
   \xrightarrow{\Delta_{2,1}}\bigwedge^2V\otimes V
   \longrightarrow S_{(2,1)}(V)\longrightarrow0
\]
is exact.  If \(\rank V=n\), the last module has rank
\(n(n^2-1)/3\).
\end{proposition}
\begin{proof}
For a one-row diagram, all exterior factors have degree one.
The relations from adjacent columns are
\(v\otimes w-w\otimes v\), so the quotient is the symmetric power.
For a one-column diagram there are no adjacent-column relations,
and the module is the indicated exterior power.

For \((2,1)\), the defining relations are precisely the image of
\(\Delta_{2,1}\).  If \(i<j<k\), that image is
\[
 (e_i\wedge e_j)\otimes e_k
 -(e_i\wedge e_k)\otimes e_j
 +(e_j\wedge e_k)\otimes e_i.
 \]
Among these three basis tensors only the last fails the
semistandard condition: the top entry \(j\) of the first column
is larger than the second-column entry \(i\).
Every nonstandard basis tensor has this unique form for one
triple \(i<j<k\).  The relation therefore eliminates that tensor
with coefficient one, and no relation for another triple
contains it.  The relations are independent and the remaining
standard tensors form a free basis of the quotient.
This proves both injectivity on the left and the asserted
basis directly.  Counting gives
\(n\binom n2-\binom n3=n(n^2-1)/3\).
\end{proof}

For \(n=3\), this computation replaces
\[
 (e_2\wedge e_3)\otimes e_1
 =-(e_1\wedge e_2)\otimes e_3
   +(e_1\wedge e_3)\otimes e_2
 \quad\hbox{in } S_{(2,1)}(V).
\]
The remaining eight tensors are a basis.  This is a concrete
straightening law, including its signs and its independence proof.
It also shows why projecting onto a preferred representation by
division by three is unnecessary.

An even smaller exact sequence is
\[
 0\longrightarrow\bigwedge^2V
  \xrightarrow{v\wedge w\mapsto v\otimes w-w\otimes v}
  V\otimes V\longrightarrow\Sym^2V\longrightarrow0.
\]
Injectivity follows because each \(e_i\wedge e_j\), \(i<j\),
has a distinct coefficient-one tensor \(e_i\otimes e_j\).
The quotient makes all tensor factors commute and has the
usual ordered monomial basis.  This proof remains valid in
characteristic two, although splitting the sequence by
symmetrization requires extra hypotheses.

\section{Rational highest weights and duals}

Let now \(k\) be algebraically closed of characteristic zero and
\(\dim V=n\).  Polynomial Schur modules indexed by partitions of
length at most \(n\) are irreducible rational representations of
\(\GL(V)\) \cite{smtABW1982}.  To describe coordinate rings of complexes one also
needs negative weights.  For a dominant integer sequence
\(\eta_1\ge\cdots\ge\eta_n\), define
\[
 S_\eta V
  =(\det V)^{\eta_n}\otimes
       S_{(\eta_1-\eta_n,\ldots,\eta_{n-1}-\eta_n,0)}V.
\]
A negative tensor power of \(\det V\) means the corresponding
power of its one-dimensional dual.

If \(\lambda,\mu\) are partitions with
\(\ell(\lambda)+\ell(\mu)\le n\), the notation
\[
 S_{(\lambda,-\mu)}V
 = S_{(\lambda_1,\ldots,\lambda_a,
           0,\ldots,0,-\mu_b,\ldots,-\mu_1)}V
\]
includes reversal of the negative block.  Without that reversal
the sequence generally would not be dominant.  In particular,
\((S_\lambda V)^*\) has highest weight
\((0,\ldots,0,-\lambda_a,\ldots,-\lambda_1)\).
The determinant twists are required even when all matrices
in the original problem have ordinary polynomial entries.

For example, if \(\dim V=2\),
\[
 S_{(q,-p)}V
   =(\det V)^{-p}\otimes\Sym^{p+q}V,
 \qquad \dim S_{(q,-p)}V=p+q+1.
\]
The dimension can be read from the monomial basis of a binary
form.  It will reappear as the dimension of a bihomogeneous
piece of the simplest nontrivial complex variety.

Over a field of positive characteristic, the integral Schur
module and its tableau basis still make sense.  Its
irreducibility and the splitting of representation filtrations
are separate assertions and generally fail.  Likewise, a
characteristic-zero Young symmetrizer can be normalized to an
idempotent only after its scalar hook factor becomes invertible.
The integral presentation above avoids making this normalization
an unspoken assumption.

\section{Schur modules applied to complexes}

A polynomial construction on modules is not automatically an
exact functor.  For the split short exact sequence
\(0\to k\to k^2\to k\to0\), applying \(\Sym^2\) term by term
gives dimensions \(1,3,1\).  Its middle homology contains the
mixed monomial and has dimension one.  Thus even splitting
the original sequence does not make this termwise sequence
exact.  A Schur complex is a construction on the entire
graded complex with its differential, not merely a list of
Schur modules of its terms.

Here is a precise construction sufficient in characteristic
zero.  If \(C^\bullet\) is a bounded complex of vector spaces,
the differential on its \(d\)-fold tensor power is
\[
 d(c_1\otimes\cdots\otimes c_d)
 =\sum_{j=1}^d
 (-1)^{|c_1|+\cdots+|c_{j-1}|}
 c_1\otimes\cdots\otimes dc_j\otimes\cdots\otimes c_d.
\]
The symmetric group permutes the homogeneous factors with
the corresponding graded signs.

\begin{proposition}[A characteristic-zero Schur complex]
Let \(d\) be a positive integer.  The displayed differential
squares to zero and commutes with
the graded permutation action.  The image of any Young
idempotent for a partition of \(d\) is a subcomplex.
If \(C^\bullet\) is acyclic, this image is acyclic.
\end{proposition}
\begin{proof}
Terms in the square that differentiate the same factor vanish
because the original differential squares to zero.  Terms
that differentiate distinct factors occur twice, with signs
differing by one degree shift; they cancel.  A transposition
of adjacent homogeneous factors commutes with the differential
by the same sign calculation.  Since adjacent transpositions
generate the symmetric group, every group algebra element
commutes with the differential, including the idempotent.
Its image and kernel are therefore complementary subcomplexes.

A bounded acyclic complex of vector spaces is contractible:
choose a complement of the cycles in each degree and invert
the differential from that complement onto the boundaries in
the next degree.  If \(h\) is the resulting contracting
homotopy, \(h\otimes\mathrm{id}\) contracts the tensor product
of \(C^\bullet\) with any other complex.  The cross terms
from the other differential cancel by the graded signs.
Thus \((C^\bullet)^{\otimes d}\) is acyclic for \(d>0\).
A direct summand of an acyclic complex has zero cohomology,
because cohomology respects a direct sum decomposition.
This proves the assertion for the idempotent image.
\end{proof}

The case \(d=0\) is deliberately excluded from the final
argument: the zeroth tensor power is \(k\), regardless of
the input complex.  Over a general coefficient ring,
divided powers and the integral Schur-complex construction
replace this idempotent argument \cite{smtABW1982}.  In either
setting, an assertion about exactness must identify the
construction actually being used.  Neither exterior powers
nor Schur terminology alone supply a new exactness theorem
for orders, a quantum algorithm, or an analytic spectral gap.

\chapter{Standard monomials on spaces of complexes}
\label{chap:standard-complexes}

\section{What a standard monomial theorem must establish}

A useful straightening theorem does three jobs.  It identifies a
set of proposed basis monomials, proves that rewriting reaches
their span, and proves that no nontrivial linear relation remains
among them.  Termination without independence is not a basis
theorem.  Independence without spanning is not a presentation of
the coordinate ring.  The examples in the next two sections
separate these tasks and prove each one.

All algebras in these sections are commutative.  The expression
\emph{standard monomial} depends on an explicitly chosen order;
it is not an intrinsic property of an arbitrary product of
determinants.  The geometry is independent of that order,
while the normal forms and the corresponding degeneration
can depend strongly on it.

\section{Rank-one matrices: the Segre example}

Let \(Z=(z_{ij})\) be an \(a\)-by-\(b\) matrix of independent
variables over a field \(k\).  Denote by \(I_2(Z)\) the ideal
of its two-by-two minors.  Consider the homomorphism
\[
 \varphi:k[z_{ij}]\longrightarrow k[u_1,\ldots,u_a,v_1,\ldots,v_b],
 \qquad z_{ij}\longmapsto u_i v_j.
\]
Every two-by-two minor belongs to its kernel.  The following
proof shows that these are all the equations, and at the
same time constructs the coordinate basis.

\begin{theorem}[Rank-one straightening, with full proof]
The kernel of \(\varphi\) is \(I_2(Z)\).  A basis of
\(k[Z]/I_2(Z)\) consists of \(1\) and the products
\[
 z_{i_1j_1}\cdots z_{i_dj_d},
 \qquad i_1\le\cdots\le i_d,\quad j_1\le\cdots\le j_d.
\]
Its degree \(d\) dimension is
\[
 \binom{a+d-1}{d}\binom{b+d-1}{d}.
\]
\end{theorem}
\begin{proof}
Use commutativity to put the row indices of any product in
nondecreasing order.  If its column sequence has an adjacent
inversion \(j>l\), the two relevant factors are
\(z_{ij}z_{kl}\) with \(i\le k\).  When \(i=k\), commuting
the two factors sorts their column indices without changing
the row sequence.  When \(i<k\), the minor relation gives
\[
 z_{ij}z_{kl}=z_{il}z_{kj}.
\]
The row sequence remains sorted, while the column inversion
count decreases by one.  Repetition terminates and proves
spanning by the displayed monomials.

Their images under \(\varphi\) are
\(u_{i_1}\cdots u_{i_d}v_{j_1}\cdots v_{j_d}\).
Distinct standard products have distinct sorted row or column
multisets, hence distinct monomial images.  Monomials in a
polynomial ring are independent, so the proposed basis is
independent.  If a polynomial lies in \(\ker\varphi\), reduce
it modulo the minor relations to a standard linear combination.
Independence forces all coefficients to be zero; this proves
the kernel assertion.  Finally, the two independently chosen
weakly increasing sequences are counted by the two displayed
binomial coefficients.
\end{proof}

The proof works without a division operation and therefore
also gives a free basis over a commutative coefficient ring.
Over a field it exhibits the coordinate ring as a subalgebra
of a polynomial ring, so it is a domain.  On the open set
\(z_{11}\ne0\), all other entries satisfy
\(z_{ij}=z_{i1}z_{1j}/z_{11}\).  The coordinates
\[
 z_{11}^{\pm1},\quad z_{21},\ldots,z_{a1},
       \quad z_{12},\ldots,z_{1b}
\]
are independent: arbitrary values with \(z_{11}\ne0\)
produce a rank-one matrix by the displayed formula.
The affine dimension is \(a+b-1\).  Its projectivization
is the Segre variety
\(\mathbb P^{a-1}\times\mathbb P^{b-1}\).

For later use, the exponent semigroup of the image is
\[
 S_{\mathrm{Segre}}
  =\{(\alpha,\beta)\in\mathbb N^a\times\mathbb N^b:
                      |\alpha|=|\beta|\}.
\]
The straightening proof explains why every such pair occurs:
list the row and column indices with those multiplicities
and pair them in sorted order.  Thus its semigroup
description is an exact description of the original ring,
not merely a limiting approximation.

\section{Two-row tableaux and the Plucker algebra}

Let \(x_1,\ldots,x_n,y_1,\ldots,y_n\) be independent
variables, and define
\[
 p_{ij}=x_i y_j-x_j y_i,\qquad 1\le i<j\le n.
\]
These are the two-by-two minors of a two-by-\(n\) matrix.
Write \(P_{ij}\) for formal generators, to distinguish a
polynomial presentation from its image in \(k[x,y]\).

\begin{lemma}[The elementary Plucker shuffle]
For \(i<k<l<j\),
\[
 p_{ij}p_{kl}=p_{il}p_{kj}-p_{ik}p_{lj}.
 \tag{P}
\]
\end{lemma}
\begin{proof}
Expand the right side using the definitions.  Terms containing
\(x_ix_j y_k y_l\) cancel, as do terms containing
\(x_kx_l y_i y_j\).  The four remaining terms are precisely
the expansion of
\((x_i y_j-x_j y_i)(x_k y_l-x_l y_k)\).
This proves the identity over \(\mathbb Z\) and hence over
any coefficient ring.
\end{proof}

\begin{theorem}[The two-row standard monomial basis]
The subalgebra \(G_{2,n}=k[p_{ij}]\) has basis
\[
 p_{i_1j_1}\cdots p_{i_dj_d},
 \qquad i_1\le\cdots\le i_d,\quad
        j_1\le\cdots\le j_d,\quad i_s<j_s.
 \tag{SM}
\]
The relations \((P)\), with \(p\) replaced by \(P\), generate
the kernel of \(k[P_{ij}]\to G_{2,n}\).
\end{theorem}
\begin{proof}
Order the factors lexicographically by their column pairs
\((i,j)\).  If the bottom row is not weakly increasing,
some two columns are incomparable for the componentwise
order.  They have indices \(i<k<l<j\), and \((P)\) replaces
their product by two products.  In each replacement the
lexicographically first changed column is either \((i,l)\)
or \((i,k)\), strictly smaller than the old \((i,j)\).
After sorting all columns again, the whole column word
has strictly decreased lexicographically.  There are
finitely many column words of the fixed degree, so
straightening terminates.  Its terminal products are
exactly \((SM)\).

For independence, order monomials lexicographically with
\(x_1>\cdots>x_n>y_1>\cdots>y_n\).
The leading monomial of \(p_{ij}\) is \(x_i y_j\).
The leading monomial of a standard product is consequently
\[
 x_{i_1}\cdots x_{i_d}y_{j_1}\cdots y_{j_d}.
\]
Such a monomial determines the two sorted rows, and those
rows determine the standard product.  Distinct standard
products therefore have distinct leading monomials.
In a nonzero linear combination, the largest of these
leading monomials cannot cancel.  This proves independence.
Reducing a polynomial in the kernel to a standard linear
combination now proves that the quadratic relations
generate the entire kernel.
\end{proof}

The tableaux have rectangular shape \((d,d)\): each column
is strictly increasing, and both rows are weakly increasing.
This gives the promised visible connection between minors
and the tableau convention of the preceding chapter.
For four indices the single relation takes the familiar
form
\[
 p_{12}p_{34}-p_{13}p_{24}+p_{14}p_{23}=0.
\]
Only the final product is nonstandard in the componentwise
column order.

The dimension can also be proved directly.  For \(n\ge2\),
localize at \(p_{12}\).  The equations imply
\[
 p_{ij}
  =\frac{p_{1i}p_{2j}-p_{1j}p_{2i}}{p_{12}},
       \qquad 3\le i<j\le n.
\]
Conversely, give arbitrary values to \(t\ne0,u_j,v_j\),
and take the matrix with columns
\[
 (1,0)^{\mathsf T},\quad(0,t)^{\mathsf T},\quad
 (-v_j/t,u_j)^{\mathsf T}\quad(j\ge3).
\]
Its indicated minors are \(p_{12}=t\), \(p_{1j}=u_j\),
\(p_{2j}=v_j\).  This proves that the localized ring is
a Laurent polynomial ring in \(2n-3\) variables.
The affine cone has dimension \(2n-3\), and the
Grassmannian \(\mathrm{Gr}(2,n)\) in Plucker coordinates
has dimension \(2(n-2)\).

\section{A complete complex example: dimensions \((1,2,1)\)}

Consider
\[
 k\xrightarrow{A} k^2\xrightarrow{B}k,\qquad
 A=\begin{pmatrix}a_1\\a_2\end{pmatrix},\quad
 B=\begin{pmatrix}b_1&b_2\end{pmatrix}.
\]
Its coordinate ring is
\[
 C=k[a_1,a_2,b_1,b_2]/(f),\qquad
 f=a_1b_1+a_2b_2.
 \tag{Q}
\]

\begin{proposition}[Geometry and basis of the basic complex]
The residue classes of monomials not divisible by \(a_2b_2\)
form a basis of \(C\).  The ring is a domain of dimension
three.  Its affine scheme is smooth away from the origin,
which is singular.  Its Hilbert series in total degree is
\[
 H_C(t)=\frac{1-t^2}{(1-t)^4}
       =\frac{1+t}{(1-t)^3},
 \qquad \dim_k C_d=(d+1)^2.
\]
With \(\deg a_i=(1,0)\) and \(\deg b_i=(0,1)\),
\(\dim_k C_{p,q}=p+q+1\).
\end{proposition}
\begin{proof}
The substitution
\[
 \begin{pmatrix}z_{11}&z_{12}\\z_{21}&z_{22}\end{pmatrix}
 =
 \begin{pmatrix}a_1&a_2\\-b_2&b_1\end{pmatrix}
\]
takes the two-by-two determinant to \(f\).
Thus the ring is the rank-one Segre ring for \(a=b=2\),
up to an invertible change of variables.  The preceding
theorem proves that it is a three-dimensional domain.

Alternatively, order the variables so that the leading
monomial of \(f\) is \(a_2b_2\).  Every multiple of \(f\)
has leading monomial divisible by \(a_2b_2\), since
leading monomials multiply in a polynomial ring over a
field.  Division by this one monic polynomial therefore
gives unique remainders supported on the stated
monomials.  This proves that particular normal form.

The gradient is
\((b_1,b_2,a_1,a_2)\), so the hypersurface is smooth at
every nonzero point by the hypersurface Jacobian
criterion.  At the origin its Zariski tangent space has
dimension four whereas the local dimension is three,
so it is singular.
Multiplication by \(f\) is injective in the polynomial
ring.  The resulting graded exact sequence gives
\(H_C(t)=(1-t^2)/(1-t)^4\), whose degree \(d\)
coefficient is \((d+1)^2\).  In bidegree \((p,q)\)
the same argument gives
\[
 (p+1)(q+1)-pq=p+q+1
\]
when \(p,q>0\); if either is zero, the identical answer
comes directly from the polynomial monomial count.
\end{proof}

Here the straightening relation is simply
\(a_2b_2=-a_1b_1\).  It is both a shuffle-type relation
and a determinant relation.  The locus with \(A\ne0\)
and \(B\ne0\) consists of complexes exact at their
middle term: the inclusion \(\im A\subseteq\ker B\)
is then an inclusion of one-dimensional subspaces.
The loci \(A=0\) and \(B=0\) are two planes in the
threefold.  Every nonzero point of either plane is
smooth.  Consequently a rank drop does not, by
itself, identify a singular point.

There is an equally explicit representation-theoretic
description.  Let the input line be \(U\), the middle
space \(V\), and the output line \(W\).  On points,
change of bases acts by
\[
 (g_U,g_V,g_W)\cdot(A,B)
 =(g_VAg_U^{-1},\,g_WBg_V^{-1}).
\]
On functions we use \((g\cdot h)(x)=h(g^{-1}x)\).
The linear coordinate spaces are therefore
\(U\otimes V^*\) and \(V\otimes W^*\).

\begin{proposition}[The bihomogeneous Schur calculation]
As a representation of this group,
\[
 C_{p,q}\cong
 U^{\otimes p}\otimes(W^*)^{\otimes q}
       \otimes(\det V)^{-p}\otimes\Sym^{p+q}V.
 \tag{QS}
\]
In characteristic zero, the middle factor is the
irreducible representation \(S_{(q,-p)}V\).
\end{proposition}
\begin{proof}
It suffices to check the middle \(\GL(V)\) action;
the powers of the outer lines follow from the
bidegrees.  The natural two-dimensional identity
\(V^*\cong(\det V)^{-1}\otimes V\), followed by
symmetric multiplication, gives a surjection
\[
 \Sym^pV^*\otimes\Sym^qV
 \longrightarrow(\det V)^{-p}\otimes\Sym^{p+q}V.
\]
In a basis with volume \(e_1\wedge e_2\), this
identity sends \(e_1^*\) to \(-e_2\) and \(e_2^*\)
to \(e_1\), with the determinant factor understood.
The contraction relation \(a_1b_1+a_2b_2\)
therefore maps to \(-e_2e_1+e_1e_2=0\).
The map factors through \(C_{p,q}\).
It is surjective because each monomial of degree
\(p+q\) can be split into factors of degrees \(p\)
and \(q\).  The source quotient and target both
have dimension \(p+q+1\), so the induced map is an
isomorphism.  The last identification follows from
the explicit highest-weight convention.
\end{proof}

This calculation checks all duals in a case where
the entire answer can be verified without a
general representation-theoretic decomposition
theorem.  Replacing \(S_{(q,-p)}V\) by
\(\Sym^pV\otimes\Sym^qV\), for example, would
give the wrong dimension and the wrong action.

\section{The general complex and its standard products}

Fix finite free modules \(V_0,\ldots,V_m\), of
ranks \(n_0,\ldots,n_m\).  We use the descending
orientation
\[
 V_m\xrightarrow{D_m}V_{m-1}\longrightarrow
       \cdots\longrightarrow V_1\xrightarrow{D_1}V_0.
 \tag{C}
\]
Thus \(D_i\) has \(n_{i-1}\) rows and \(n_i\)
columns, and the equations are \(D_iD_{i+1}=0\).
At the endpoints we use \(D_0:V_0\to0\) and
\(D_{m+1}:0\to V_m\), both zero.
Let \(A_R\) be the polynomial ring in all entries
modulo all entries of these products.

For \(I\subseteq[n_{i-1}]\) and
\(J\subseteq[n_i]\) with \(|I|=|J|=p>0\),
write \([I\mid J]_i\) for the corresponding minor.
Within a single matrix, order minors by
\[
 [I\mid J]_i\preceq[I'\mid J']_i
 \quad\Longleftrightarrow\quad
 p\ge q,\quad i_t\le i'_t,\quad j_t\le j'_t
          \quad(1\le t\le q),
 \tag{M}
\]
where \(q=|I'|\).  Minors of different matrices
are declared incomparable in this partial order.
This order by itself is not enough to describe
standard products across different arrows.

Here is the additional rule, stated in the
transpose convention of De Concini--Eisenbud--Procesi,
Section 16 \cite{v2DEP1982}.  For a pair of minors,
interchange the two factors if necessary so that
their arrow indices satisfy \(i\ge j\).
The pair is standard if one of the following holds:

\begin{enumerate}
\item \(i\ge j+2\);
\item \(i=j\) and the two minors are comparable
      under \((M)\);
\item \(j=i-1\), and the following complement
      condition holds.
\end{enumerate}

For the last condition write
\(x=[I\mid J]_i\), of size \(p\), and
\(x'=[I'\mid J']_{i-1}\), of size \(q\).
The shared vector space has dimension
\(n=n_{i-1}\).  List
\([n]\setminus I=\{\ell_1<\cdots<\ell_{n-p}\}\)
and \(J'=\{j'_1<\cdots<j'_q\}\).  The condition is
\[
 q\le n-p,\qquad
 j'_t\le \ell_{\,n-p-q+t}\quad(1\le t\le q).
 \tag{MC}
\]
A product of minors is standard when every pair
of its factors is standard.  For \((Q)\), the
rule says that \(a_2b_2\) is nonstandard while
the other three mixed pairs are standard,
exactly as the direct calculation found.

\begin{theorem}[External standard basis theorem for complexes]
The standard products just defined form an \(R\)-basis
of \(A_R\).  For nonnegative integer bounds \(r_i\), if all minors of
\(D_i\) larger than \(r_i\) are set equal to zero, the surviving standard
products form a basis of the quotient.  These
constructions commute with base change.
Over a field, the quotients are reduced.
If \(r_0=r_{m+1}=0\) and
\[
 r_i+r_{i+1}\le n_i\qquad(0\le i\le m),
 \tag{A}
\]
the rank-bounded quotient is a normal
Cohen--Macaulay domain.
\end{theorem}

The basis and reducedness assertions are
De Concini--Eisenbud--Procesi, Theorem 16.1 and
Corollary 16.2; the final assertion follows from
their Theorem 16.3 and Proposition 16.4,
which in turn refer to De Concini--Strickland
\cite{v2DEP1982,v2DS1981}.
The general topological and commutative-algebraic
arguments establishing these results are not
replaced here by the elementary examples.
In particular, the existence of a resolution
of a variety would not alone prove its normality.

Two features of the statement deserve attention.
First, the basis includes products of minors
of different sizes, whose polynomial degrees
are those sizes.  It is not a basis just in
the degree-one matrix entries.  Second, since
nonadjacent arrows may form standard products
despite their minors being incomparable in
\((M)\), this is a generalized Hodge algebra.
Its standardness is not simply the chain
condition in one ordinal partial order.
That distinction will determine which
simplicial complex appears in a monomial
degeneration.

\section{Rank strata and orbit closures without a source shortcut}

For the following geometric argument let \(k\) be
algebraically closed.  Define \(C_r\) to be the
closed locus in \((C)\) with
\(\rank D_i\le r_i\), where \(r\) satisfies
\((A)\).  Let \(C_r^\circ\) denote the locus
with all these ranks equal to \(r_i\).
The use of equality here is essential:
strictly smaller ranks would not define
the orbit indexed by \(r\).

\begin{proposition}[Orbits, closure, and dimension]
The locus \(C_r^\circ\) is a single orbit under
\(\prod_{i=0}^m\GL(V_i)\), and
\(C_r=\overline{C_r^\circ}\).
It is irreducible for every admissible \(r\).
The irreducible components of the full space
of complexes are precisely the \(C_r\)
for coordinatewise maximal admissible \(r\).
Moreover
\[
 \dim C_r
  =\sum_{i=1}^m r_i(n_i+n_{i-1}-r_i)
        -\sum_{i=1}^{m-1}r_i r_{i+1}.
 \tag{D}
\]
\end{proposition}
\begin{proof}
At each vertex choose subspaces
\[
 V_i=B_i\oplus H_i\oplus T_i,\qquad
 B_i=\im D_{i+1},\quad B_i\oplus H_i=\ker D_i.
\]
The dimensions are \(r_{i+1}\),
\(n_i-r_i-r_{i+1}\), and \(r_i\),
respectively.  The only nonzero block of \(D_i\)
is an isomorphism \(T_i\to B_{i-1}\).
Changing bases makes all these isomorphisms
identity matrices.  Thus the ranks determine
a unique orbit.

Decrease one rank by one by replacing a
single identity entry in the corresponding
block by a parameter \(t\).  For \(t\ne0\)
the old ranks remain; at \(t=0\) precisely
that rank drops.  The compositions stay zero
because the direct-sum blocks have not changed.
Iterating shows that every smaller rank orbit
is in the closure.  Conversely, the vanishing
of the \((r_i+1)\)-minors is closed and prevents
any larger rank from entering the closure.
This proves the closure assertion.
The group is irreducible, hence so is its
orbit closure.  There are finitely many
admissible rank vectors; their containment
order is coordinatewise.  The maximal ones
therefore give exactly the components.

For the dimension, first choose flags
\(B_i\subset Z_i\subset V_i\), with
\(\dim B_i=r_{i+1}\) and \(\dim Z_i=n_i-r_i\).
Choosing \(B_i\) contributes
\(r_{i+1}(n_i-r_{i+1})\) parameters.
Choosing \(Z_i/B_i\) in \(V_i/B_i\)
contributes
\(r_i(n_i-r_i-r_{i+1})\).
Finally choose each isomorphism
\(V_i/Z_i\to B_{i-1}\), an open subset
of a space of dimension \(r_i^2\).
Every complex in the stratum determines
these flags and isomorphisms uniquely.
Adding the dimensions and cancelling
the repeated squares gives \((D)\).
\end{proof}

This proof also corrects a possible ambiguity in
phrases such as \emph{the irreducible rank sequence}.
Every admissible rank bound gives an irreducible
orbit closure; maximality is needed for being
an irreducible component of the full space.
The equations imposing rank at most \(r_i\)
use minors of size \(r_i+1\), not \(r_i\).
These distinctions are visible in the normal
form itself and do not depend on terminology.

\section{The Schur decomposition with all vertices and duals}

Return to an algebraically closed field \(k\)
of characteristic zero.  On the descending
complex, let
\[
 (g_0,\ldots,g_m)\cdot D_i
       =g_{i-1}D_i g_i^{-1},
 \qquad (g\cdot f)(D)=f(g^{-1}\cdot D).
 \tag{G}
\]
The linear coordinates on the \(i\)th arrow
are therefore \(V_i\otimes V_{i-1}^*\).
For a tuple of partitions
\(\lambda=(\lambda_1,\ldots,\lambda_m)\),
put \(\lambda_0=\lambda_{m+1}=\varnothing\).
Require \(\ell(\lambda_i)\le r_i\).
At vertex \(j\), define the dominant weight
\[
 \eta_j(\lambda)=
 \bigl(\lambda_j,\,
       0^{\,n_j-\ell(\lambda_j)-\ell(\lambda_{j+1})},\,
       -\operatorname{rev}(\lambda_{j+1})\bigr).
 \tag{W}
\]
Condition \((A)\) ensures that the number
of zero entries is nonnegative.

\begin{theorem}[External Schur filtration and its characteristic-zero splitting]
The coordinate ring \(k[C_r]\) has a multiplicative
tableau filtration whose successive modules are
\[
 \bigotimes_{j=0}^m S_{\eta_j(\lambda)}V_j,
 \qquad \ell(\lambda_i)\le r_i.
 \tag{SF}
\]
The summand indexed by \(\lambda\) has arrow
multidegree \((|\lambda_1|,\ldots,|\lambda_m|)\).
In characteristic zero this gives a
\(\prod_j\GL(V_j)\)-module isomorphism
\[
 k[C_r]\cong
 \bigoplus_{\substack{\lambda_1,\ldots,\lambda_m\\
                     \ell(\lambda_i)\le r_i}}
          \bigotimes_{j=0}^m S_{\eta_j(\lambda)}V_j.
 \tag{SD}
\]
\end{theorem}

The filtration is the Schur straightening
construction of De Concini--Strickland,
reviewed explicitly in Carroll--Weyman,
Section 4 \cite{v2DS1981,smtCW2011}.
Their tableaux concatenate an incoming
dual diagram, expressed using a determinant
twist, with the outgoing diagram.
Formula \((W)\) is that construction in
the action convention \((G)\).
Rank bounds remove exactly the partitions
whose column length exceeds the allowed
rank.  Each finite multidegree is a
rational representation of a reductive
group in characteristic zero, so its
filtration splits as a representation.
This last step does not assert a
canonical splitting as an algebra.

There are \(m+1\) factors in \((SF)\),
including both endpoints.  Omitting the
last vertex, or treating both appearances
of an intermediate vector space as
covariant, gives a wrong representation.
For one arrow the formula reduces to
the Cauchy decomposition
\[
 \Sym^d(V_1\otimes V_0^*)
   \cong \bigoplus_{|\lambda|=d}
                   S_\lambda V_1\otimes S_\lambda V_0^*,
\]
with the partition length cut off by a
rank bound when one is present.
For \((Q)\), take \(\lambda_2=(p)\)
and \(\lambda_1=(q)\).  The middle
weight is exactly \((q,-p)\), and
\((SD)\) reproduces the independently
proved formula \((QS)\).

Multiplication contains further information.
Multiplying two tableau expressions of
shapes \(\lambda,\mu\) lies in the filtration
piece indexed by \(\lambda+\mu\); straightening
can add terms in lower pieces.  The
associated graded multiplication retains
the top piece, often described in terms
of highest-weight or Cartan multiplication.
One cannot reconstruct every product in
the original ring simply by writing a
direct sum of representation spaces.
An isomorphism of modules answers a
linear question; an isomorphism of
algebras must also answer this
multiplicative question.

\chapter{Associated graded rings, semigroups, and toric limits}
\label{chap:filtered-toric}

\section{Three different constructions}

There are three constructions in the preceding discussion
that should be distinguished.  A Schur filtration organizes
the coordinate ring into representations.  A straightening
degeneration replaces nonstandard products by zero and
produces a monomial quotient.  A suitable subalgebra
degeneration replaces determinants by their diagonal
monomials and can produce an integral semigroup ring.
They can start from the same ring and end in different
rings.  The difference is mathematical, not just a
choice of notation.

For an ascending, exhaustive filtration
\(F_0A\subseteq F_1A\subseteq\cdots\) with
\(1\in F_0A\) and \(F_iA\,F_jA\subseteq F_{i+j}A\),
the associated graded algebra is
\[
 \operatorname{gr}_F A
   =\bigoplus_{d\ge0}F_dA/F_{d-1}A,\qquad F_{-1}A=0.
\]
Its product keeps only the highest filtration class
of a product.  A nonzero product in \(A\) can fall
into a lower filtration piece and therefore become
zero in the associated graded ring.

The Rees algebra records the lost information:
\[
 \mathcal R_F(A)=\bigoplus_{d\ge0}F_dA\,t^d
       \subseteq A[t].
\]
The quotient at \(t=1\) recovers \(A\), while the
quotient by \(t\) recovers \(\operatorname{gr}_F A\).
For the latter statement, the degree \(d\)
coefficient of \(t\mathcal R_F(A)\) is
\(F_{d-1}A\), which gives exactly the graded
quotient above.  The former follows by evaluation
at one and finite telescoping of the coefficients.
Over a field, choose a vector-space basis
adapted to the filtration by successively
extending a basis of \(F_{d-1}A\) to one
of \(F_dA\).  If a chosen basis vector
\(b\) first appears at stage \(\delta(b)\),
the elements \(b\,t^{\delta(b)}\) form
a \(k[t]\)-basis of the Rees algebra:
every later occurrence of \(b\) is a
multiple by a power of \(t\), and
independence follows inside \(A[t]\).
Thus this algebra is free, hence flat,
over \(k[t]\).  Finite generation as
an algebra requires a suitable finite
filtration; it is not automatic for
every abstract filtration.

Flatness preserves the Hilbert function in the
graded situations used below.  It does not
assert that every fiber is irreducible, smooth,
or a domain.  The following explicit
construction makes that warning verifiable.

\section{A monomial degeneration for a quadratic straightening algebra}

Let \(P\) be a finite partially ordered set,
and let \(A\) be a graded algebra generated
by elements \(x_p\), \(p\in P\), all of
degree one.  Suppose that the chain products
\[
 x_{p_1}\cdots x_{p_d},\qquad
                p_1\le\cdots\le p_d,
\]
form a basis.  Suppose also that every
incomparable pair \(p,q\) has a straightening
relation
\[
 x_px_q=\sum_{u\le v}c_{uv}x_ux_v,
 \qquad c_{uv}\ne0\Longrightarrow u<p,\ u<q.
 \tag{ASL}
\]
These are the homogeneous quadratic hypotheses
needed for the next argument.  General Hodge
algebras can have generators of different
degrees, so they must not be substituted
into this statement without checking those
hypotheses.

\begin{proposition}[An explicit straightening degeneration]
Under these assumptions, \(A\) has a flat
graded family over \(k[t]\) whose fiber at
one is \(A\) and whose fiber at zero is
\[
 k[X_p:p\in P]/
        (X_pX_q:p,q\hbox{ incomparable}).
 \tag{SR}
\]
The chain monomials form a free \(k[t]\)-basis
of the family.
\end{proposition}
\begin{proof}
If \(P\) is empty, then \(A=k\) and the family \(k[t]\)
proves the assertion. Assume from now on that \(P\ne\varnothing\).
Let \(h(p)\) be the maximum number of strict
steps in a chain starting at \(p\) and going
upward.  Put \(D_p=3^{h(p)}\), choose an integer
\(M>\max D_p\), and set \(w_p=M-D_p>0\).
If a term \(x_ux_v\) occurs on the right
of \((ASL)\), then
\[
 D_u\ge3\max(D_p,D_q)>D_p+D_q.
\]
It follows that \(w_u+w_v<w_p+w_q\).
Thus straightening strictly decreases the
positive integer weight of a monomial.

In a polynomial presentation, every element
reduces by these relations to a combination
of chain monomials.  Since their images
are independent in \(A\), a relation reduces
to zero.  This proves both that the
straightening relations generate the
presentation ideal and that their initial
ideal for this weight order is exactly
the ideal of incomparable products.
In particular no additional chain monomial
is killed by taking initial terms.

Now replace each relation by
\[
 X_pX_q-\sum c_{uv}
       t^{\,w_p+w_q-w_u-w_v}X_uX_v.
 \]
All displayed exponents of \(t\) are
strictly positive.  The same decreasing
weight argument gives spanning over
\(k[t]\) by chain monomials.  After
inverting \(t\), the change
\(X_p=t^{w_p}x_p\) identifies this
algebra with \(A[t,t^{-1}]\).
If a \(k[t]\)-linear combination of
chain monomials were zero, this change
would give a relation among the
corresponding basis monomials of \(A\),
with Laurent polynomial coefficients.
Each coefficient would be zero.
The original polynomial coefficients
are therefore zero as well.
The family is free over \(k[t]\), hence
flat, and specializing its relations
gives the two stated fibers.
\end{proof}

The quotient \((SR)\) is the
Stanley--Reisner ring of the order
complex of \(P\): a subset is a face
exactly when it is a chain.  Whenever
there are incomparable \(p,q\), the
images of \(X_p\) and \(X_q\) are
nonzero but their product is zero.
Thus this ring has zero divisors.

For a two-by-two rank-one matrix,
write its entries \(p,q,r,s\), with
relation \(ps-qr=0\).  The only
incomparable pair in the grid order
is \(q,r\).  The monomial limit is
\[
 k[p,q,r,s]/(qr),
\]
the union of the two coordinate
hyperplanes \(q=0\) and \(r=0\).
The original ring is a domain by
the Segre proof.  Calling this
monomial limit a toric domain would
therefore be an error.

For varieties of complexes, the
general Hodge-algebra construction
gives a stepwise Rees degeneration
to the ring governed by its
nonstandard products
\cite[Section 3]{v2DEP1982}.
Because the forbidden products
are squarefree pairs, that
discrete ring is the
Stanley--Reisner ring of a flag
simplicial complex: a set is
a face exactly when all its
pairs are standard.  It need
not be the order complex of
the minor order \((M)\), since
standardness across different
arrows uses the extra rule
\((MC)\).

This identifies precisely what
the combinatorial basis supplies.
It supplies a free module basis,
a multiplication algorithm, and
a controlled degeneration.
Normality and the Cohen--Macaulay
property of the original rank
variety come from the stronger
theorems cited in the preceding
chapter.  They cannot be inferred
merely by observing that some
associated graded ring has
monomial equations.

\section{Initial subalgebras and the diagonal semigroup}

An initial subalgebra is a different
construction from an initial
quotient.  If \(G\subseteq k[z_1,\ldots,z_N]\)
is a subalgebra and a term order
has been chosen, define
\[
 \operatorname{in}(G)
   =k[\operatorname{LM}(f):0\ne f\in G],
\]
where \(\operatorname{LM}\) denotes
the leading monomial.  A set of
generators is called a SAGBI basis
if their leading monomials
generate this entire initial
subalgebra.  This is the
subalgebra counterpart of a
Groebner basis; arbitrary
algebra generators do not
automatically have this property.

\begin{proposition}[The two-row minors are a SAGBI basis]
For the lexicographic order
\(x_1>\cdots>x_n>y_1>\cdots>y_n\),
the minors \(p_{ij}\), \(i<j\),
are a SAGBI basis of \(G_{2,n}\).
Their initial subalgebra is
\[
 k[x_i y_j:1\le i<j\le n].
 \tag{T}
\]
\end{proposition}
\begin{proof}
Expand any nonzero element of
\(G_{2,n}\) in the standard
monomial basis \((SM)\).
Those basis elements have
distinct leading monomials,
as proved earlier.  The
largest of their leading
monomials cannot cancel.
It is a product of the
\(x_i y_j\).  Thus every
leading monomial of an
element of \(G_{2,n}\)
belongs to the subalgebra
on the right of \((T)\).
The reverse inclusion holds
because each \(x_i y_j\)
is the leading monomial
of the corresponding
minor.  This proves the
claim.
\end{proof}

Unlike \((SR)\), the ring
\((T)\) is a subring of a
polynomial ring and is
therefore a domain.  Its
presentation uses binomial
equalities between products
having the same exponent
vector.  Those products do
not become zero.

For clarity, an affine
semigroup here means a
finitely generated additive
submonoid \(S\subseteq
\mathbb N^N\).  Its
semigroup ring has basis
symbols \(z^s\), \(s\in S\),
and multiplication
\(z^s z^{s'}=z^{s+s'}\).
The embedding in a
polynomial ring is part
of this definition.
For \((T)\), the generators
of \(S\) are the vectors
\((e_i,e_j)\), \(i<j\),
in \(\mathbb N^n\times
\mathbb N^n\).

\begin{theorem}[The diagonal semigroup and its inequalities]
This semigroup consists exactly of the pairs
\((\alpha,\beta)\in\mathbb N^n\times\mathbb N^n\)
satisfying
\[
 |\alpha|=|\beta|,\qquad
 \sum_{j=1}^k\beta_j
       \le\sum_{i=1}^{k-1}\alpha_i
       \quad(1\le k\le n).
 \tag{I}
\]
Each such pair determines a unique
semistandard two-row tableau by
sorting its top and bottom rows.
The semigroup is saturated in
the lattice it generates.
\end{theorem}
\begin{proof}
An expression as a sum of
generators is a matching of
top indices \(i\) with
bottom indices \(j>i\),
with multiplicities.  Any
bottom index at most \(k\)
must be matched to a top
index smaller than \(k\).
This proves all the prefix
inequalities and the equality
of total sizes.

Conversely, process the bottom
indices in order \(j=1,\ldots,n\).
Before processing \(j\),
the available top-index
capacity with \(i<j\) is
\[
 \sum_{i<j}\alpha_i-\sum_{h<j}\beta_h.
 \]
All previous matches used
indices in that set.
Inequality \((I)\) says
that the available capacity
is at least \(\beta_j\).
Choose that many units of
capacity, making edges
from the available \(i<j\)
to \(j\).  The process
never fails.  At the end,
the equality of totals
shows that all capacity
has been used.  Thus the
pair is a sum of the
specified generators.

Sort the top indices as
\(i_1\le\cdots\le i_d\)
and the bottom indices
as \(j_1\le\cdots\le j_d\).
If \(j_s=k\), the number
of bottom entries at
most \(k\) is at least
\(s\).  The prefix
inequality ensures that
at least \(s\) top
entries are smaller
than \(k\).  Hence
\(i_s<j_s\), so these
rows form a semistandard
tableau.  Their
multiplicities uniquely
determine the sorted
rows, proving uniqueness.

Finally let \(L\) be
the subgroup generated
by \(S\), and suppose
\(qv\in S\) for an
integer \(q>0\) and
\(v\in L\).  Every
coordinate of \(v\)
is nonnegative because
the same is true of
\(qv\).  The equality
and inequalities in
\((I)\) are homogeneous,
so they also hold for
\(v\).  The characterization
just proved gives
\(v\in S\).
\end{proof}

The first and last prefix
conditions force
\(\beta_1=0\) and
\(\alpha_n=0\), as
they should: there is
no column with bottom
entry one or top
entry \(n\).  These
boundary conditions
are sometimes lost
when a tableau
semigroup is described
only by equal total
degrees.

\section{Gelfand--Tsetlin patterns in the example}

Put
\[
 A_r=\sum_{i\le r}\alpha_i,\qquad
 B_r=\sum_{j\le r}\beta_j,\qquad
 A_0=B_0=0.
\]
The conditions \((I)\) become
\[
 A_r\ge A_{r-1}\ge B_r\ge B_{r-1},
 \qquad A_n=B_n=d.
 \tag{GT}
\]
The pair \((A_r,B_r)\)
is the shape of the
subtableau formed by
entries at most \(r\).
Its interlacing with
the preceding pair
is exactly the
two-row Gelfand--Tsetlin
condition.  Taking
successive differences
recovers \(\alpha,\beta\),
so \((GT)\) is an
equivalent integral
description, not a
second source of
unverified basis
vectors.

Addition of exponent
vectors becomes
componentwise addition
of the patterns.
This is the semigroup
operation in the
toric limit.  In the
original determinant
algebra, multiplying
the two corresponding
standard products
may require
straightening, with
several terms.  The
pattern addition
therefore records
the leading product,
not every coefficient
in the original
multiplication table.

The saturation in the
preceding theorem also
has a direct
commutative-algebraic
consequence.

\begin{proposition}[Normality from these integral inequalities]
The semigroup rings
\(k[S_{\mathrm{Segre}}]\)
and \(k[S]\) of \((T)\)
are normal domains.
\end{proposition}
\begin{proof}
Both semigroups are
the intersection of
their generated lattice
\(L\) with finitely
many homogeneous
integral linear
half-spaces: coordinate
nonnegativity and,
for \(S\), the prefix
inequalities.  The
required total-degree
equalities hold on
the lattice \(L\);
the inequalities
select its
semigroup points.
This follows from
their explicit
characterizations.

The group ring \(k[L]\)
is a Laurent
polynomial ring
after choosing a
basis of the free
abelian group \(L\).
It is a normal
domain.  Its
fraction field is
also the fraction
field of \(k[S]\),
because every lattice
element is a
difference of two
semigroup elements.
For each defining
integer linear
functional \(\ell\),
give a nonzero
Laurent polynomial
the value
\[
 v_\ell\left(\sum c_u z^u\right)
       =\min_{c_u\ne0}\ell(u).
 \]
The lowest homogeneous
parts of two Laurent
polynomials have
nonzero product, so
these values add
under multiplication.
They extend by
subtraction to a
valuation of the
fraction field.

Suppose \(f\) belongs
to this common
fraction field and
is integral over \(k[S]\).
It is integral over
\(k[L]\), so belongs
to that Laurent
polynomial ring.
It belongs to
every valuation
ring \(v_\ell\ge0\):
otherwise an integral
equation for \(f\)
would have its
monic highest power
as the unique term
of smallest valuation,
which cannot sum
to zero.  Thus all
exponents of \(f\)
satisfy every
defining inequality.
They lie in \(S\),
and \(f\in k[S]\).
This proves integral
closedness.  The
domain assertion
follows from the
monomial embeddings.
\end{proof}

This proof uses the
actual inequality
description; it
does not infer
normality from
finite generation
alone.  A finitely
generated semigroup
can fail saturation,
as \(\langle2,3\rangle
\subseteq\mathbb N\)
shows: \(2\cdot1\)
belongs to it but
\(1\) does not.
Correspondingly,
\(k[t^2,t^3]\) is
not normal, since
\(t\) is integral
over it and lies
in its fraction
field without
belonging to the
ring.

\section{Two flat families from the same Grassmannian}

The difference between
the two degenerations
can be made entirely
explicit for \(n=4\).
Let \(P\) denote its
six Plucker variables
and form the two
\(k[t]\)-algebras
\[
 \begin{split}
 \mathcal A_{\mathrm{mon}}
 &=k[t,P]/
   (P_{14}P_{23}-tP_{13}P_{24}+tP_{12}P_{34}),\\
 \mathcal A_{\mathrm{tor}}
 &=k[t,P]/
   (P_{14}P_{23}-P_{13}P_{24}+tP_{12}P_{34}).
 \end{split}
 \tag{F}
\]
Both fibers at \(t=1\)
are the affine
Plucker cone over
\(\mathrm{Gr}(2,4)\).

\begin{proposition}[Flatness and different special fibers]
Both algebras in \((F)\)
are free over \(k[t]\),
with basis the
monomials not divisible
by \(P_{14}P_{23}\).
Their special fibers
are, respectively,
\[
 k[P]/(P_{14}P_{23}),\qquad
 k[P]/(P_{14}P_{23}-P_{13}P_{24}).
 \]
The first has zero
divisors and the
second is a normal
domain.
\end{proposition}
\begin{proof}
Choose a term order
on the six \(P\)
variables for which
\(P_{14}P_{23}\)
is larger than
the other two
products.  Each
relation is monic
in this leading
monomial.  Division
over the coefficient
ring \(k[t]\)
gives a remainder
supported on the
listed monomials.
The leading term
of a nonzero
multiple of the
relation is
divisible by
\(P_{14}P_{23}\),
so no nonzero
remainder lies
in the ideal.
This proves
uniqueness and
freeness.

At zero, the
first relation
is the product
of two distinct
variables, which
proves the
zero-divisor
assertion.
The second
relation is a
two-by-two
determinant
in four of
the variables.
Its ring is
the two-by-two
Segre ring
with two
independent
polynomial
variables adjoined.
The preceding
normality proof,
or the same
semigroup
argument with
two free
coordinates,
proves it is
a normal
domain.
\end{proof}

For any nonzero
value of \(t\),
rescaling \(P_{14}\)
in the first
family and
\(P_{12}\) in
the second
identifies the
fiber with the
original quadric.
Thus these are
actual one-parameter
degenerations,
not merely two
unrelated examples.
In both cases
the degree \(d\)
dimension is
\[
 \binom{d+5}{5}-\binom{d+3}{5}
   =\frac{(d+1)(d+2)^2(d+3)}{12},
\]
where the second
binomial is zero
when \(d<2\).
Agreement of this
Hilbert function
does not make
the two special
fibers isomorphic.
One is reducible
and the other
is integral.

The general study
of toric limits
of flag and
Schubert varieties
extends this
pattern; it
requires its
own filtration
or subalgebra
argument
\cite{smtGL1997}.
The completely
proved two-row
case here is
sufficient to
show why an
ordinary
straightening
degeneration
cannot be
renamed a
toric degeneration
without that
additional
argument.

\section{How the algebra enters the larger research program}

The construction gives
a precise interface
with invariant theory.
At a vertex \(j\)
of the descending
complex, a scalar
matrix \(c_jI\)
acts on the
Schur factor
\((W)\) by
\[
 c_j^{\,|\lambda_j|-|\lambda_{j+1}|}.
 \]
This follows
either by summing
the highest-weight
entries or
directly from
the source and
target coordinate
actions in
\((G)\).
It is therefore
possible to
test necessary
central-character
conditions for
an invariant
without expanding
every polynomial.
To obtain a
semi-invariant
one must also
identify the
appropriate
one-dimensional
determinant
character at
each vertex.
Balancing scalar
degrees alone
does not prove
that an entire
Schur module is
one-dimensional.

For gentle algebras
whose colored
paths describe
complexes, the
same vertex
space can appear
in several such
factors.  Taking
invariants then
requires decomposing
their tensor
product, or
applying the
precise
semi-invariant
criterion.
Carroll--Weyman
carry out that
additional step
under their
quiver hypotheses
\cite{smtCW2011}.
The following
chapters give
the matching
semigroup and
determinantal
examples.  The
coordinate ring
of a single
complex and
the semi-invariant
ring of a
gentle algebra
are distinct
objects, even
when both
admit useful
tableau
descriptions.

There is also a
useful integral
interface with
the theory of
orders.  The
exterior and
shuffle formulas
are integral,
and the standard
bases are free
over the base
ring.  One can
therefore perform
the same finite
normal-form
calculations
before and
after reduction
modulo a prime.
This controls
the ranks of
the indicated
graded modules
and their
Hilbert
functions.
The
characteristic-zero
decomposition
into irreducible
representations
does not survive
merely because
the basis does;
semisimplicity
is an additional
property of
the coefficient
field.

For computation,
the bases give
exact finite
representations
of functions
on the rank
varieties.
The straightening
algorithm can
multiply those
functions and
reduce equations
to normal form.
The leading
semigroup can
organize enumeration
and Hilbert
function checks.
These are
algebraic
algorithms with
costs that depend
on degrees,
dimensions, and
coefficient
sizes.  A basis
theorem supplies
no automatic
polynomial-time
algorithm for
elliptic-curve
discrete
logarithms and
no error
analysis for
a lattice
cryptosystem.
Such claims
would need
their own
reduction and
complexity
analysis.

The same separation
applies to
spectral questions.
Schur characters,
Hilbert functions,
and \(K\)-theory
classes record
algebraic or
representation-theoretic
information.
They do not
specify a
Hilbert-space
inner product,
a self-adjoint
Hamiltonian, or
a uniform lower
bound on its
positive spectrum.
A flat family
can preserve
every graded
dimension while
changing its
singularities.
A family of
finite exact
complexes can
likewise have
positive
eigenvalues
tending to
zero.  The
mass-gap
transport
problem
therefore
requires the
analytic
hypotheses
stated elsewhere
in this
monograph.

\reviewstatus{The exterior basis, shuffle identities,
contraction homotopy, elementary Schur sequences,
characteristic-zero tensor-complex argument,
Segre and two-row Plucker bases, basic-quadric
calculations, rank-orbit dimension formula,
quadratic straightening family, diagonal
semigroup description, its normality, and
the two explicit flat families are proved
in these chapters.  The general Schur
basis theorem, general complex standard
basis and normality theorems, and general
Schur filtration are cited external
theorems with stated hypotheses.  The
Lean files verify selected interfaces
described in the formalization appendix;
they do not formalize these general
standard monomial theorems.}



